\documentclass[11pt]{article}
\usepackage[letterpaper,margin=0.9in]{geometry}
\usepackage{amsmath,amssymb}
\usepackage[hidelinks]{hyperref}
\newcommand{\F}{\mathbb F_q}
\newcommand{\ket}[1]{\lvert #1\rangle}
\newcommand{\bra}[1]{\langle #1\rvert}
\newcommand{\ip}[2]{\langle #1,#2\rangle}
\newcommand{\one}{\mathbf 1}
\newcommand{\file}[1]{{\small\nolinkurl{#1}}}
\setlength{\parskip}{4pt}
\setlength{\parindent}{0pt}
\title{Reading a hidden hyperplane from a coherent sample}
\author{}\date{}
\begin{document}
\maketitle
\vspace{-3.2em}
Let $q$ be prime, $r\ge2$, and $f\in\F^r$ an unknown nonzero vector. Given
copies of the state
\[
  \ket{H_f}=q^{-(r-1)/2}\sum_{x:\,f\cdot x=1}\ket{x},
\]
we want to read all $r$ entries of $f$. At $q=17$ and $r=64$, the optimal
one-copy probability is $0.907558$, and two copies give $0.996494$.
One classical uniform point on the same plane leaves exactly $17^{63}$
equally likely candidates under the uniform prior. We construct a circuit
attaining the quantum optimum for every copy count $k\ge1$. Its other use is
as a coherent decoder for quantum information stored in superpositions
of parallel planes.

\section*{1. The operator in the Lean development}
In the archive, lines 11--29 of
\file{Formalism/LogVandermondeBounds/FinrankFiberSumZero.lean}
define the incidence operators $U,D$ with entries
$A_{x,f}=\one_{\{f\cdot x=1\}}$, the sum $\mathrm{Nm}$ along each punctured
scalar line, and the total sum $\mathrm{tr}$ over nonzero points. We use
the Euclidean inner product and omit the zero row and column. Then $D=A$,
$U=A^*$, and the column of $D/q^{(r-1)/2}$ at $f$ is $\ket{H_f}$.

In \file{LogVandermondeProofs/RamifiedLatticeOperatorIdentity.lean},
\file{count1} and \file{count2} (lines 180 and 192) count the solutions of one
and two independent equations. The latter proof finds an invertible
two-coordinate minor. The proof of \file{DU_formula} (line 228) splits the
double incidence sum according to equality, proportionality and independence.
It obtains, for $x\ne0$,
\[
  (DUh)(x)=q^{r-1}h(x)-q^{r-2}\mathrm{Nm}\,h(x)+q^{r-2}\mathrm{tr}\,h .
\]
On amplitudes summing to zero along each punctured line, the last two terms
vanish. The transposed identity gives the same conclusion on the column side.

The archive also gives the operator a Fourier interpretation. With
$\psi(t)=e^{2\pi i t/q}$ let $Gh(f)=\sum_{c\ne0}\psi(c)h(cf)$.
The file \file{FourierIsometryOnNormKernel.lean} proves
$GG^*=qI-\mathrm{Nm}$ at line 84. Its second theorem,
\file{fourier_eq_down_comp_gauss} (115), gives the unnormalized
additive Fourier transform $\mathcal F=DG$ on the line-sum-zero subspace.
The sum over $f\cdot x=0$ vanishes line by line.
In \file{FiniteModelOfAdmissibleTransports.lean}, \file{Unf.DU} (277) and
\file{Unf.FF_eq} (384) transfer the two identities through the local
coordinate charts. The proof calls \file{FF_DGs} (529) twice inside
\file{FF_iso} (555), then iterates that inner-product identity in
\file{FFlist} (568). This file is in the import closure of \file{src/Main.lean}.
These are source-level observations. The algorithm and its optimality proof
below use finite linear algebra independently of the status of GRH.

\newpage
\section*{2. The optimum from three intersection counts}
Choose the representative $d$ of each projective direction whose first
nonzero entry is $1$. There are $m=(q^r-1)/(q-1)$ directions. Write
$H_{d,b}=\{x:d\cdot x=b\}$. We consider $B=\F^\times$, which gives exactly
the Lean family via $f=d/b$, and $B=\F$, which includes planes through the
origin. Put $s=|B|$. The input is $k$ identical coherent copies for a single
hidden $(d,b)$. The output is a classical label. The source supplies copies
directly, with no membership, preparation or inverse-preparation queries.

For $d=e$, distinct offsets $b,c$ give disjoint supports. For $d\ne e$,
the two planes share $q^{r-2}$ points, giving inner product $1/q$.
Tensoring $k$ times raises each inner product to the $k$th power. Thus the
Gram matrix on the $ms$ labels, with $J_n$ the all-ones matrix, is
\[
  G_B= I_{ms}+q^{-k}(J_m-I_m)\otimes J_s.
\]
Write $a=s/q^k$. Its spectrum, with multiplicities, is
\[
\begin{array}{lcc}
\text{label amplitudes}&\text{eigenvalue}&\text{multiplicity}\\ \hline
\text{sum zero across offsets within every direction}&1&m(s-1)\\
\text{constant within directions, overall sum zero}&1-a&m-1\\
\text{constant on all labels}&1+(m-1)a&1 .
\end{array}
\]
In particular the normalized incidence matrix preserves an entire
$m(q-2)$-dimensional space of coherent signals for $B=\F^\times$ and $k=1$.
It preserves their inner products even when the signals superpose different
directions.

\textbf{Optimal measurement.} For either $B$, the maximum uniform-prior
success probability, also the maximum worst-case probability, is
\[
 \boxed{\displaystyle P_k=
 \left(\frac{s-1+(1-1/m)\sqrt{1-a}+m^{-1}\sqrt{1+(m-1)a}}{s}\right)^2 .}
\]
To prove this, collect the $k$-copy input columns in $V$, so $G_B=V^*V$.
Every diagonal entry of $\sqrt{G_B}$ is
$h=\operatorname{Tr}(\sqrt{G_B})/(ms)$ by the three spectral projections.
The polar isometry of $V$ has an adjoint that sends the input columns to
the columns of $\sqrt{G_B}$. Reading the label gives success $h^2$ on
every input. This is the square-root measurement. For a direct upper
bound set $R=(VV^*)^{1/2}$ on its support. Each column $v_j$ satisfies
$v_j^*R^{-1}v_j=h$ by singular value decomposition. Cauchy--Schwarz
gives $v_jv_j^*\le hR$. For any POVM $\{E_j\}$ the mean success is at most
$\operatorname{Tr}(hR)/(ms)=h^2$. All inverses here act on the support,
including the singular case $k=1$, $s=q$. This recovers the symmetric
pure-state discrimination principle in this particular geometry \cite{eldar}.

For classical uniform points from the Lean family, $t\ge1$ observations
form the rows of a matrix $X$. Under a uniform prior the posterior is
uniform on the $q^{r-\mathrm{rank}(X)}$ solutions of $Xf=\one$.
For $t<r$ the success probability is at most $q^{t-r}$, regardless of
postprocessing. The comparison counts coherent samples and classical
random points from their respective sources.

\newpage
\section*{3. A circuit that attains the bound}
The ordinary Fourier measurement reveals $d$ on a nonzero readout.
Its outcome probabilities are independent of $b$. To retain $b$ we keep
the scalar frequency coherent. The Fourier transform with positive phase
on each point register gives
\[
  q^{-1/2}\sum_{t\in\F}\psi(bt)\ket{td}
\]
per copy. First allow all offsets. For $k=1$, relabel each nonzero vector
$td$ as $(d,t)$ and map the shared zero vector to
$m^{-1/2}\sum_e\ket{e,0}$. The output labels for nonzero $t$ are
orthogonal to this vector, so this is an isometry. Apply the inverse
one-dimensional Fourier transform to $t$ and measure. The correct label
has amplitude $(q-1+1/\sqrt m)/q$ for every plane, including $b=0$.

For $k\ge1$, combine the $k$ Fourier registers as follows.
\begin{enumerate}
\item For a nonzero tuple $(t_1d,\ldots,t_kd)$, extract $d$ from its first
nonzero vector. Replace the scalars by $u=\sum_j t_j$ and
$(t_2,\ldots,t_k)$. For $u\ne0$ the latter register is uniform on all
$q^{k-1}$ tuples. For $u=0$ and a nonzero input tuple it is uniform on the
$q^{k-1}-1$ nonzero tuples. Undo these uniform-state preparations in
their respective branches. Keep the all-zero input as a separate symbol
$\ket{*}$. When $k=1$ the second uniform set is empty and its branch
has amplitude zero.
\item Put $\theta=q^{1-k}$ and $\ket{\bar d}=m^{-1/2}\sum_e\ket e$.
Up to the common factor $1/\sqrt q$, the $u=0$ component now equals
$\sqrt{1-\theta}\ket{d,0}+\sqrt{\theta}\ket{*}$. In the plane
spanned by $\ket{\bar d,0}$ and $\ket{*}$ rotate
\[
  \frac{\sqrt{1-\theta}\ket{\bar d,0}+\sqrt{m\theta}\ket{*}}
       {\sqrt{1+(m-1)\theta}}
   \quad\longmapsto\quad \ket{\bar d,0}.
\]
Act by the identity on the orthogonal complement of this plane. The unused
output in this plane has zero amplitude for every promised input.
\item Invert the scalar Fourier transform on $u$. For the full offset set
this already implements the optimal decoder. Indeed, in the Fourier basis
on offsets it acts on the input label map by amplitude gain $1$ at $u\ne0$.
At $u=0$ the gains are $\sqrt{1-\theta}$ on differences of directions and
$\sqrt{1+(m-1)\theta}$ on the uniform direction. Thus the map from
labels through encoding and decoding is precisely $\sqrt{G_{\F}}$.
\end{enumerate}

\textbf{Taking advantage of the known nonzero offset.}
When $B=\F^\times$, the final correction acts only on offset means. Split
the direction register coherently into its uniform subspace and its
orthogonal complement. On the offset register in each sector put
$T=I+(\sqrt\lambda-1)\ket{+}\bra{+}$, where $\ket{+}$ is the uniform
$q$-offset state and $\lambda$ is respectively $1+(m-1)\theta$ or
$1-\theta$. Step 3 implemented $T$ in that sector. Let
$\ket B=s^{-1/2}\sum_{b\in B}\ket b$. Rotate the real unit vector
$T\ket B/\|T\ket B\|$ to $\ket B$ within
$\mathrm{span}\{\ket B,\ket0\}$, fixing its orthogonal complement.
Here $\|T\ket B\|^2=1+(\lambda-1)s/q>0$.
Every difference between offsets in $B$ is fixed. Hence on the allowed
inputs this correction produces exactly $\sqrt{G_B}$, with zero amplitude
on disallowed labels. Reading $(d,b)$ and returning $d/b$ recovers $f$.

\newpage
\section*{4. Implementation, coherent signals and checks}
\textbf{Resources.} The circuit uses $k$ copies on $kr$ field registers.
Projective normalization uses a nonzero-coordinate test and one field
inversion per recovered direction. Directions with first nonzero coordinate
at $j$ form a block of size $q^{r-j}$ for $j=1,\ldots,r$. These contiguous
blocks give efficient reversible ranking and unranking, hence preparation
of $\ket{\bar d}$ from a uniform integer in $[0,m)$. The other uniform
states use intervals of lengths $q$, $q-1$, $q^{k-1}$ and $q^{k-1}-1$.
Uniform interval states can be prepared one bit at a time with amplitudes
determined by the counts in each child interval. Unpreparing a uniform
state reduces a rotation about it to a controlled low-dimensional rotation.
The map on tuples extends reversibly by retaining the tuple coordinates
until their uniform state is unprepared. The promised inputs have no
tuples spanning two different directions.

With reversible integer arithmetic, approximate Fourier transforms over
$\mathbb Z/q$ and synthesized rotations, the gate count is polynomial in
$k,r,\log q,\log(1/\epsilon)$ for overall error $\epsilon$. For example a
cyclic Fourier transform can use phase estimation on the modular shift
with extra precision, followed by coherent rounding of the frequency
and unpreparation of the Fourier eigenstate from its known frequency.
Repeated squaring of the shift amounts to modular addition by known
constants. The probabilities in this note refer to ideal gates. Coherent
state preparation from a classically stored list would have its own cost.

\textbf{Decoding a quantum payload.} Let $c_{d,b}$ be arbitrary complex
amplitudes with $\sum_{b\ne0}c_{d,b}=0$ for each $d$ and
$\sum_{d,b}|c_{d,b}|^2=1$. Applying the one-copy decoder to
\[
     \sum_{d,b\ne0}c_{d,b}\ket{H_{d,b}}
     \quad\text{produces}\quad
     \sum_{d,b\ne0}c_{d,b}\ket{d,b}.
\]
The zero-frequency amplitude cancels inside each direction. The first
register can therefore identify a single occupied direction while the
second retains an unknown $(q-2)$-dimensional quantum state. Across
multiple directions the entire linear isometry also preserves entanglement
with any reference system. The inverse circuit encodes these signals.
This is the line-sum-zero condition used by the arithmetic proof, with
nonzero offsets corresponding to the nonzero scalar multiples of a normal.

\textbf{Independent numerical check.} \file{hyperplane_decoder.py} builds
each plane in the computational basis, tensors copies, applies the additive
Fourier transform numerically, compresses uniform fibers, and carries out
the rotations above. It checks the full complex amplitude matrix against
the positive square root obtained by an independent singular value
decomposition of the input columns. Eleven choices of $(q,r,k)$ over
$q=2,3,5,7$ test both offset families and up to three copies. The maximum
entrywise discrepancy across the checks was below $4\times10^{-15}$.
The $r=64$ example is an evaluation of the exact formula.

\textbf{Related algorithms.} Fourier sampling for hidden flats and
quantum ridgelet transforms already reorganize spatial frequencies by
direction \cite{childs,ridge}. Here the calculation gives the coherent
treatment of the shared zero frequency and the optimal probability of
recovering direction and offset together for the stated copy model.
Priority for this specialization has not been established.

\small
\begin{thebibliography}{9}
\setlength{\itemsep}{0pt}\setlength{\parskip}{0pt}
\bibitem{eldar} Y. C. Eldar, A. Megretski and G. C. Verghese,
\href{https://arxiv.org/abs/quant-ph/0211111}{Optimal detection of symmetric mixed quantum states}.
\bibitem{childs} A. M. Childs, L. J. Schulman and U. V. Vazirani,
\href{https://users.cms.caltech.edu/~schulman/Papers/quantum-nonlinear.pdf}{Quantum algorithms for hidden nonlinear structures}.
\bibitem{ridge} H. Yamasaki, S. Subramanian, S. Hayakawa and S. Sonoda,
\href{https://arxiv.org/abs/2301.11936}{Quantum ridgelet transform}.
\end{thebibliography}
\end{document}
